\documentclass[11pt]{article}

\input{0-parameters.tex}

\begin{document}


\section{Using costs numbers from the CCC}

\subsection{Mathematical framework}

Recall that there exists a value of $\zeta$  such that we can get aggregate productivity from
\begin{equation}
\hat{A}_{t,\tau} = \sum_i \frac{p_{i,t}X_{i,t}}{\text{GDP}_{t}} \left(\zeta \hat{X}_{i,t,\tau} + (\zeta - 1) \hat{p}_{i,t,\tau} \right)
\label{eq:main_result_recall}
\end{equation}

To fix ideas, imagine that both steel and cars becomes more expensive. To get a correct prediction for aggregate productivity, we need to include price increases of both industries into this formula, even if the car industry was the only buyer of steel. 


Here we do not have price growth, but many institutions, including the CCC, attempt to compute ``total extra costs'' from the transition, say in billions of pounds. Here we explain how these total costs numbers can be fit into this formula.

Let $C_t$ denote the annualised extra cost of the transition at time $t$. At some future time $t +\tau$, the \textit{additional} annualised extra cost is $\Delta C_{i,t,\tau} = C_{i,t+\tau} - C_{i,t}$, so we can write the (unit) price growth between $t$ and $t+\tau$ as 
\begin{equation}
\hat{p}_{i,t,\tau} = \frac{\Delta C_{i,t,\tau}/X_{i,t}}{p_{i,t}},
\label{eq:phat_from_C}    
\end{equation}
Using \eqref{eq:phat_from_C} and the fact that the ``extra cost'' is derived under the assumption of constant services ($\hat{X}_{i,t,\tau}=0$), Eq. \ref{eq:main_result_recall} nicely simplifies to
   
\begin{equation}
 \hat{A}_{t,\tau} = (\zeta - 1)  \frac{\sum_i \Delta C_{i,t,\tau}}{\text{GDP}}.
\label{eq_ahat_t_tau_ccc}
\end{equation}

This equation is remarkable in that it makes it clear that it does not matter in which industry the extra costs occur. We only need to count total extra incurred costs. However, we do need to count the extra costs in \textit{all} industries. For instance, if both steel and cars become more expensive, we need to add the extra costs of both industries, even if cars were more expensive only because of the extra cost of steel.

To get some intuition for this, consider electricity, and imagine that it has a productivity shock. Aggregate productivity will be affected as $\hat{A} = \lambda_{\text{elec}} \hat{A}_{\text{elec}}$ \--- when we know productivity shocks, we do not need to know what happens to other sectors (recalling we assume that productivity shocks are independent). However, if we do not know productivity shocks, but only prices and quantities, we can still recover aggregate productivity, but \textit{we do need to know prices and quantity change for all industries}, not only in the electricity industry. Indeed, if electricity has a productivity shock and becomes cheaper, many other industries will have cheaper prices (even though their TFP is constant by assumption).

This clarifies how our framework deals with supply chain interdependencies. When we can access productivity shocks, we can derive aggregate impacts  without taking into the supply chain \--- Domar aggregation is a ``network irrelevance result''. When we cannot access productivity shocks, but only price and quantity growth, then according to Eq. \ref{eq:main_result_recall}, we can recover aggregate productivity growth but only if we know the correct value of $\zeta$ and the prices and quantity growth in all industries, even if only a few industries have non-zero productivity shocks.


\subsection{Unprocessed and processed costs}

The CCC offers raw, ``unprocessed costs'', and ``processed'', free of double counting costs.

To understand the difference, consider the process of costing the transition. A straightforward approach is to list all the industries concerned by the transition, and for each industry, evaluate how much it costs to provide the same goods/services before and after the transition. For instance, it could be that we will need to pay more for steel, and more for cars.

One can argue that to evaluate the total costs of the transition, we cannot simply sum up the extra costs of each industry. This is because the cost of an industry depends on the costs of other industries, which can lead to double counting of the costs. In the example above, the extra cost of steel is taken into account when computing the cost of cars; consumers do not need to pay the extra cost of steel twice (once to the steel industry and once to the car industry), but only once, to the car industry.

Thus, the CCC offers ``free of double counting'' costs, where it adds together the extra cost of steel, the extra costs of cars, but deduct the extra cost of steel already included in the cost of cars. Is this right for us? While we cannot know exactly how the CCC evaluates and removes double counting, we think it essentially evaluates, for each industry, the part of cost increases that are due to cost increases of inputs. Let us define it as
\[
\Delta C_{i} = \underbrace{\Delta C^U_{i}}_{\text{cost increase already incurred upstream}} +  \underbrace{\Delta C^*_{i}}_{\text{residual cost}}
\]

Now, starting from the dual definition of TFP
\begin{align}
\label{eq_dual_processed}
\hat{A}_i &= \hat{p}_i - \sum_j a_{ji} \hat{p}_{j} \\
\hat{A}_i &= \frac{1}{p_i X_i} \left( (p_i X_i) \hat{p}_i - (p_i X_i) \sum_j a_{ji} \hat{p}_{j} \right)\\
\hat{A}_i &= \frac{1}{p_i X_i} \left( (p_i X_i)\hat{p}_i - \sum_j X_{ji} \hat{p}_{j}\right)\\
\hat{A}_i &= \frac{1}{p_i X_i} \left( \Delta C_i  - \Delta C_i^{U} \right) = \frac{\Delta C^*_i}{p_i X_i},
\end{align}
where the definition of $\Delta C_i$ comes from Eq. \ref{eq:phat_from_C}, and the identity
\[
\Delta C_i^{U} = \sum_j X_{ji} \hat{p}_{j}
\]
is our interpretation of what the CCC counts as ``double counted'', that is, total costs incurred only due to extra costs of the inputs. The last line makes clear that ``free of double counting'' costs essentially correspond to TFP shocks.

Indeed, inserting this into Domar aggregation,
\begin{align}
\label{eq_ahat_processed_ccc}
\hat{A} &= \sum_i \frac{p_i X_i}{\text{GDP}} \hat{A}_i  = \frac{\sum_i \Delta C^*_i}{\text{GDP}}.
\end{align}

Comparing this to Eq. \ref{eq_ahat_t_tau_ccc}, we see that if we use the costs ``free of double counting'', we do not need to pre-multiply by $(\zeta-1)$. Essentially, this is because ``free of double counting'' costs correspond to TFP shocks, while ``raw'' costs correspond to price changes.


\section{Comments Emilien 14th April}
\subsection{Correcting typo}

There's a slight mistake in Eq. \eqref{eq_dual_processed}, see below

\begin{align}
\hat{A}_i &= \er{-(}\hat{p}_i - \sum_j a_{ji} \hat{p}_{j}\er{)} \\
\hat{A}_i &= \er{-(}\frac{1}{p_i X_i} \left( (p_i X_i) \hat{p}_i - (p_i X_i) \sum_j a_{ji} \hat{p}_{j} \right)\er{)}\\
\hat{A}_i &= \er{-(}\frac{1}{p_i X_i} \left( (p_i X_i)\hat{p}_i - \sum_j X_{ji}\er{p_{j}} \hat{p}_{j}\right)\er{)}\\
\hat{A}_i &= \er{-(}\frac{1}{p_i X_i} \left( (p_i X_i)\frac{\Delta p_i}{p_i} - \sum_j X_{ji}p_{j} \frac{\Delta p_j}{p_j} \right)\er{)}\\
\hat{A}_i &= \er{-(}\frac{1}{p_i X_i} \left( X_i\Delta p_i - \sum_j X_{ji} \Delta p_j \right)\er{)}\\
\hat{A}_i &= \er{-(}\frac{1}{p_i X_i} \left( \Delta C_i  - \Delta C_i^{U} \right)\er{)} = \er{-(}\frac{\Delta C^*_i}{p_i X_i}\er{)},
\end{align}

Correcting it, we get
\begin{equation}
    \hat{A}_i  = -\frac{\Delta C^*_i}{p_i X_i}.
\end{equation}


\subsection{Back to square one?}

François, you've proposed to focus the paper on the methodology and the CCC scenarios. In this application, the main equation we use is $\hat{A}  = -\frac{\sum_i \Delta C^*_i}{GDP}$. Meaning that the aggregate productivity growth impact is the ratio of the extra net cost in 2050 and 2022 GDP. We find that the impact is -0.6\%, because the extra cost are 0.6\% of GDP.

It is actually pretty similar to the DICE-like approach. In \cite{dietzEndogenousGrowthConvexity2015}, GDP is 
\[
Y_t=(1-D_t^Y)(1-\Lambda_t)A_tK^{\alpha+\beta}_tL^{1-\alpha}_t
\]. $D_t$ is the damage to the economy, and $\Lambda_t$ is the emissions abatement costs as a share of GDP. $A_t$ is the exogenous element of TFP. One can express the endogenous TFP as $(1-\Lambda_t)A_t$. If we assume there is no growth of the exogenous element (or it is outside the scope of the analysis. Then $\hat{A}_t=\ln(1-\Lambda_t)\approx \Lambda_t= \frac{\sum_i \Delta C_{i,t,\tau}}{\text{GDP}}$. The impact of net zero on aggregate productivity is the sum of the net cost over GDP.

We are then back to what we were doing and did not satisfy you two years ago. And most importantly, the whole framework is not useful in deriving Eq. \eqref{eq_ahat_processed_ccc}. 

\paragraph{My take} We get back to this classic (and intuitive) result, only because we are implicitly making the assumption that $\zeta=0$ (see below). What we are adding is that if we don't make that simplifying assumption, then it just need a correction factor in front of it to take into account the dual nature of TFP (basically, the influence of $\hat{X}_i$).



\subsection{Calibrating $\zeta$ on the CCC scenario}

You show that by writing the dual we get for each sector that ``free of double counting'' costs correspond to TFP shocks, while ``raw'' costs correspond to price changes. Equating Eq. \eqref{eq_dual_processed} and Eq. \eqref{eq_ahat_t_tau_ccc}, we get
\begin{equation*}
    -\Delta C^*_i=(\zeta-1)\Delta C_i,
\end{equation*} 
hence 
\begin{equation}
\label{eq_zeta_ccc}
     \zeta=\frac{-\Delta C^*_i}{\Delta C_i}+1
\end{equation}

Let us look at that in the CCC data. In Table \ref{tab_TFP_annualised_nonprocessed_processed}, for each sector is the annualised resource cost (annualised processed cost) given by the CCC. My best effort at recomputing it (given that I don't have the discount rate nor lifespan). Last column is the annualised unprocessed cost using the same discount rate and lifespan as the annualised processed cost. The zeta columns correponds to Eq. \eqref{eq_zeta_ccc}.


\begin{longtable}{cP{3cm}P{3cm}P{3cm}P{3cm}}
\caption{Extra 2050 cost (bn£) per sector before/after
the re-allocation to remove double counting (Balanced Net
Zero, CCC) and compared to the CCC annualised resource cost - Annualised cost approach. \footnotesize \textit{Manual correction of the Fuel supply as the model does work great to annualise it for the non-processed cost. To be investigated}} \\ 
  \hline
Sector & CCC annual cost & Recomputed annual cost (Processed) & Recomputed annual cost (Non-processed) & zeta \\ 
  \hline
Agriculture & -0.99 & 0.18 & 0.19 & 0.08 \\ 
  Aviation & -1.27 & -0.93 & -0.86 & -0.08 \\ 
  Electricity supply & -5.30 & -2.62 & 1.40 & 2.88 \\ 
  F-gases & -0.00 & 0.01 & 0.01 & 0.00 \\ 
  Fuel supply & -1.56 & -1.69 & - & - \\ 
  LULUCF sinks & 2.62 & 2.53 & 2.53 & 0.00 \\ 
  LULUCF sources & -0.09 & 0.05 & 0.05 & 0.00 \\ 
  M\&C & 4.18 & 4.11 & 5.02 & 0.18 \\ 
  Non res buildings & 3.81 & 3.86 & 3.68 & -0.05 \\ 
  Removals & 5.71 & 6.36 & 6.57 & 0.03 \\ 
  Res buildings & 12.90 & 11.57 & 15.64 & 0.26 \\ 
  Shipping & 3.62 & 3.58 & 3.44 & -0.04 \\ 
  Surface transport & -9.09 & -7.79 & 1.29 & 7.03 \\ 
  Waste & 1.69 & 1.75 & 3.49 & 0.50 \\ 
   \hline
Total & 16.21 & 20.96 & 42.45 & 0.81 \\ 
  \hline
\label{tab_TFP_annualised_nonprocessed_processed}
\end{longtable}

Using the aggregate cost, we can compute productivity growth
\[
\hat{A}=-\frac{\sum_i \Delta C^*_i}{\text{GDP}}=-\frac{16.21}{2234}=-0.007
\]

With my own approximation of annualisation
\[
\hat{A}=-\frac{\sum_i \Delta C^*_i}{\text{GDP}}=-\frac{20.1}{2234}=-0.009
\]

With the non-processed costs, using  Eq. \eqref{eq_ahat_t_tau_ccc} and the calibrated $\zeta=0.13$
\[
\hat{A}=(\zeta-1)\frac{\sum_i \Delta C_i}{\text{GDP}}=(0.13-1)\frac{42.45}{2234}=-0.017
\]

We should have a $\zeta=\frac{-\sum_i \Delta C^*_i}{\sum_i \Delta C_i}+1=-\frac{20.1}{42.45}+1=0.52$ to account for the gap between processed and un-processed costs. 



\subsection{Do we just assume that $\zeta=0$ in the CCC application?}

We show a paradox when forcing $\hat{X}_M=0$ implying that $\zeta=0$.

\vspace{3mm}

As a reminder we consider this the simplified economy
\begin{itemize}
    \item \textbf{Upstream sector $U$} using only labour. 
    \[
    \hat{X}_U > 0,~
    \hat{p}_U < 0
    \]
    \item \textbf{Downstream sector $M$} using $U$ as its sole input. 
    \[
    \hat{X}_M = \hat{A}_M+a_{UM}  \hat{X}_U,~
    \hat{p}_M = -\hat{A}_M+a_{UM}  \hat{p}_U
    \]
\end{itemize}

% As soon as we assume that $\hat{X}_M=0$ and $\hat{X}_U=0$, then we get a problem. From the primal we get $\hat{X}_M = \hat{A}_M+a_{UM}  \hat{X}_U$, therefore, $\hat{A}_M=0$.

% It might not be a problem in a more complicated economy because we can enforce $\forall i, ~\hat{X}_i=0$ but still have $\hat{X}_{ji}\neq0$. Thus a reallocation of the production between sectors. 

We have two sectors, aggregate productivity is
\begin{equation*}
    \hat{A}=\lambda_M \hat{A}_M+\lambda_U \hat{A}_U.
\end{equation*}

As shown before, we consider only the downstream sector M, and show that we retrieve the true aggregate productivity growth.

\begin{equation*}
    \begin{aligned}
    \hat{\tilde{A}}=\lambda_M  \hat{\tilde{A}}_M 
    &= \lambda_M\Bigl(\zeta \hat{X}_M + (\zeta-1) \hat{p}_M\Bigr)\\
    &= \lambda_M\Bigl(\zeta(\hat{A}_M+ \tfrac{\lambda_U}{\lambda_M} \hat{X}_U )
        +  (\zeta-1) (-\hat{A}_M+\tfrac{\lambda_U}{\lambda_M} \hat{p}_U)\Bigr)\\
    &= \lambda_M (\zeta\hat{A}_M-(\zeta-1)\hat{A}_M)+\lambda_U\Bigl(\zeta \hat{X}_U + (\zeta-1) \hat{p}_U\Bigr)\\
    &= \lambda_M \hat{A}_M+\lambda_U \hat{\tilde{A}}_U\\
    &=\lambda_M \hat{A}_M+\lambda_U \hat{A}_U\quad(\text{since we calibrate }\zeta\text{ such that }\hat{\tilde{A}}_U=\hat{A}_U)\\
    &=\hat{A}.
\end{aligned}
\end{equation*}

Hence, equating the first and the last lines, we get
\begin{equation}
\label{eq_hatA_framework_1}
 \hat{A}=    \lambda_M\Bigl(\zeta \hat{X}_M + (\zeta-1) \hat{p}_M\Bigr).
\end{equation}

Which means, that in the particular case $\hat{X}_M=0$, we get $ \hat{A}=    \lambda_M(\zeta-1) \hat{p}_M.
$

But if I directly take into account that $\hat{X}_M=0$, we get
\begin{equation}
\label{eq_hatA_framework_1_5}
    \begin{aligned}
    \hat{\tilde{A}}=\lambda_M  \hat{\tilde{A}}_M 
    &= \lambda_M\Bigl((\zeta-1) \hat{p}_M\Bigr)\\
    &= \lambda_M(\zeta-1) (-\hat{A}_M+\tfrac{\lambda_U}{\lambda_M} \hat{p}_U)\\
    &= \lambda_M(1-\zeta) \hat{A}_M +\lambda_U(\zeta-1)\hat{p}_U
\end{aligned}
\end{equation}

Note that from the dual equation that $\hat{A}_U=-\hat{p}_U$, hence
\begin{equation} 
\label{eq_hatA_framework_2}
     \hat{\tilde{A}}= (1-\zeta)\left(\lambda_M \hat{A}_M +\lambda_U\hat{A}_U\right)=
     (1-\zeta)\hat{A}.
\end{equation}

These two approaches, equations Eq. \eqref{eq_hatA_framework_1} and Eq. \eqref{eq_hatA_framework_2}, should give the same result. They do if and only if $\zeta=0$. 

We make that assumption for the upstream sector, equation the equation from the dual  $\hat{A}_U=-\hat{p}_U$, and the equation from our framework,  $\hat{A}_U=(\zeta-1)\hat{p}_U$, which imposes $\zeta=0$ for U. 

If we do not make that assumption, from Eq. \eqref{eq_hatA_framework_1_5}, we get
\begin{equation*}
    \hat{\tilde{A}}=\lambda_M(1-\zeta) \hat{A}_M +\lambda_U\hat{A}_U,
\end{equation*}
which is not equal to the true productivity growth.
And especially, for the $U$ sector, we impose $\hat{A}=(\zeta-1)\hat{p}_U$ with $\zeta\neq0$ which contradicts the dual equation which states that dual  $\hat{A}_U=-\hat{p}_U$.

\paragraph{Conclusion} We are making an error in all cases: if we assume that $\zeta=0$ then we have a trivial model (see previous section), if not we are making an error on computing the productivity of the upstream sector. 






\newpage
\section{Appendix: Avoiding double counting}

We are considering in the CCC application (and in later applications to some key sectors) only the key sectors for the net zero. 
But there is a risk that these sectors see decrease in price, not because they're more productive, but because their inputs are being produced more productively.
In a simplified vision:  An upstream industry $U$ and a downstream sector $M$ consuming $U$ as an input. If $M$ appears more productive (i.e. offers a lower price) only because it purchases cheaper inputs from $U$, do we need to model $U$ explicitly to capture the net productivity gains? 


\paragraph{Toy Model\\}


Let the economy be composed of two sectors:
\begin{itemize}
    \item An upstream industry $U$ (e.g., producing solar PV modules, batteries).
    \item A downstream sector $M$ (e.g., the electricity generation sector that purchases solar PV from $U$).
\end{itemize}

Consider the simplified economy described, such as
\begin{itemize}
    \item \textbf{Upstream sector $U$} using only labour. 
    \[
    \begin{cases}
    \hat{X}_U > 0,\\
    \hat{p}_U < 0,\\
    0 < \zeta < 1.
    \end{cases}
    \]
    \item \textbf{Downstream sector $M$} using $U$ as its sole input. Let $a_{UM}$ be the cost share of $U$ in $M$'s expenses.  By using primal and dual expressions in productivity growth accounting we can express
    \[
    \begin{cases}
    \hat{X}_M = \hat{A}_M+a_{UM}  \hat{X}_U,\\
    \hat{p}_M = -\hat{A}_M+a_{UM}  \hat{p}_U,\\
    \hat{A}_M > 0.
    \end{cases}
    \]
\end{itemize}

\paragraph{Relating Cost Shares and Domar Weights}
The share of $U$ in $M$'s total expenses is
\begin{equation*}
     a_{UM}  =  \frac{p_U X_U}{p_M X_M}  =  \frac{p_U X_U}{p_Y Y} \frac{p_Y Y}{p_M X_M} 
     =  \frac{\lambda_U}{\lambda_M},
\end{equation*}
where $\lambda_U = \tfrac{p_U X_U}{Y}$ and $\lambda_M = \tfrac{p_M X_M}{Y}$ are the Domar weights of $U$ and $M$, with $p_Y Y$ is total nominal GDP of the economy. Hence, 
\begin{equation*}
    \hat{X}_M = \hat{A}_M+\frac{\lambda_U}{\lambda_M} \hat{X}_U
    \quad\text{and}\quad
    \hat{p}_M = -\hat{A}_M+\frac{\lambda_U}{\lambda_M} \hat{p}_U.
\end{equation*}

\paragraph{True Aggregate Productivity.}
Aggregate productivity growth $\hat{A}$ is 
\begin{equation}
      \hat{A} 
     =  
    \lambda_U  \hat{A}_U  +  \lambda_M  \hat{A}_M.
\end{equation}


\paragraph{Downstream-Only Productivity Measure} We now show that a simplified measure of productivity that focuses solely on the downstream sector $M$ still recovers $\hat{A}$. In our framework, we define a downstream ``effective'' productivity growth $\hat{\tilde{A}}_M$ as
\begin{equation*}
    \hat{\tilde{A}}_M =  \zeta \hat{X}_M + (\zeta - 1) \hat{p}_M.
\end{equation*}
Then the overall productivity, from the perspective of $M$ alone, is scaled by the Domar weight $\lambda_M$:
\begin{equation}
    \begin{aligned}
    \hat{\tilde{A}}=\lambda_M  \hat{\tilde{A}}_M 
    &= \lambda_M\Bigl(\zeta \hat{X}_M + (\zeta-1) \hat{p}_M\Bigr)\\
    &= \lambda_M\Bigl(\zeta(\hat{A}_M+ \tfrac{\lambda_U}{\lambda_M} \hat{X}_U )
        +  (\zeta-1) (-\hat{A}_M+\tfrac{\lambda_U}{\lambda_M} \hat{p}_U)\Bigr)\\
    &= \lambda_M (\zeta\hat{A}_M-(\zeta-1)\hat{A}_M)+\lambda_U\Bigl(\zeta \hat{X}_U + (\zeta-1) \hat{p}_U\Bigr)\\
    &= \lambda_M \hat{A}_M+\lambda_U \hat{\tilde{A}}_U\\
    &=\lambda_M \hat{A}_M+\lambda_U \hat{A}_U\quad(\text{since we calibrate }\zeta\text{ such that }\hat{\tilde{A}}_U=\hat{A}_U).
\end{aligned}
\end{equation}
We know that $\hat{\tilde{A}}_U = \hat{A}_U$ because we calibrate $\zeta$ to exactly match observed productivity growth.
This confirms that our method correctly captures the full productivity dynamics by considering only the downstream industry. The reason it works relies on two key points: (i) in standard growth accounting, sector $M$'s productivity growth does not reflect the price decline of sector $U$, whereas here we account for it as part of $M$'s productivity; (ii) using $\lambda_M a_{UM} = \lambda_U$, we ensure the correct Domar weight appears in the decomposition. 

Hence,
\begin{equation}
      \hat{A}  =  \lambda_U \hat{A}_U+\lambda_M \hat{A}_M 
     =  \lambda_M \hat{\tilde{A}}_M.
\end{equation}
In short, analysing the downstream sector alone (with a suitable proxy for productivity growth) captures the true aggregate productivity growth.

\paragraph{What if we use the two sectors?}
If we use the price increase and quantities increase of both U and M, then we compute
\begin{align*}
 \hat{\tilde{A}}=\lambda_M  \hat{\tilde{A}}_M +  \lambda_U  \hat{\tilde{A}}_U 
    &=
    \lambda_M  \hat{A}_M +  2\lambda_U  \hat{A}_U.
\end{align*}

If we only consider U, then we predict
\begin{equation*}
 \hat{\tilde{A}}=  \lambda_U  \hat{\tilde{A}}_U= \lambda_U  \hat{A}_U.
\end{equation*}

Let us assume we consider either U, M or both U and M to apply how framework to compute our prediction of aggregate productivity growth $\hat{\tilde{A}}$. In which case are we making the biggest error $\varepsilon=\hat{A}-\hat{\tilde{A}}$?

\begin{equation}
\begin{cases}
    U:~ \hat{A} \approx  \lambda_U  \hat{A}_U, ~\varepsilon= -\lambda_M  \hat{A}_M,\\
      M:~ \hat{A} \approx  \lambda_U \hat{A}_U+\lambda_M \hat{A}_M , ~\varepsilon= 0,\\
        U\&M:~ \hat{A} \approx  \lambda_M  \hat{A}_M +  2\lambda_U  \hat{A}_U , ~\varepsilon= \lambda_U  \hat{A}_U.\\
\end{cases}
\end{equation}

If we consider the two sectors, then we are making a relative error on aggregate productivity 
\begin{equation*}
   \frac{\varepsilon}{\hat{A}}= \frac{\lambda_U  \hat{A}_U}{ \lambda_U \hat{A}_U+\lambda_M \hat{A}_M}<<1
\end{equation*}
which is small if $\lambda_U<<\lambda_M$, which is likely to be the case and unlikely to be compensated by the magnitude of the relative productivity shock.

Hence, the conclusion is, it's better to consider the ``downstream'' sector which has both input changes and idiosyncratic productivity shock if we have prior information on it. If not, then the error we're making is still relatively small if one sector is smaller than the other.

\subsubsection{Correcting for double counting}

In the CCC case, where they give us the extra OPEX and CAPEX, we transform them into price growth. They make sure to exclude any double counting by focusing on “non-financed costs” -- costs not matched by corresponding payments. For instance, increased electricity investment due to EV adoption is financed by higher household electricity spending, so only one side of the transaction should be counted. 

As in section \ref{section_results_ccc_method}, let $C_t$ denote the annualised extra cost of the transition at time $t$. The \textit{additional} annualised extra cost between $t$ and $t+\tau$ is $\Delta C_{i,t,\tau} = C_{i,t+\tau} - C_{i,t}$, the (unit) price growth between $t$ and $t+\tau$ is $\hat{p}_{i,t,\tau} = \frac{\sfrac{\Delta C_{i,t,\tau}}{Q_i}}{p_{i,t}}$, where $Q_i$ is the number of units, and $p_i$ the price.
Assuming quantity growth is null, we get in our framework 
\begin{equation}
 \hat{A}_{i,t,\tau} = (\zeta - 1) \frac{\Delta C_{i,t,\tau}}{p_{i,t} Q_{i}}.
\end{equation}

\begin{equation}
 \hat{A}_{t,\tau} = \sum_i \frac{p_{i,t} Q_{i}}{\text{GDP}} \Big( (\zeta - 1) \frac{\Delta C_{i,t,\tau}}{p_{i,t} Q_{i}} \Big) = (\zeta - 1)  \frac{\sum_i \Delta C_{i,t,\tau}}{\text{GDP}}.
\end{equation}


We consider two sectors, $U$ and $M$ as in the previous section. They undergo extra cost $\Delta C_{U,t,\tau}$ and $\Delta C_{M,t,\tau}=\Delta C_{U,t,\tau}+\Delta \tilde{C}_{M,t,\tau}$, \FL{where the first term represents extra bills paid to $U$, and the second term represents the cost increases of $M$ not due to increasing expenses in $U$.}

\begin{equation}
\begin{cases}
    \hat{A}_{M,t,\tau} = (\zeta - 1) \frac{\Delta C_{M,t,\tau}}{p_{M,t} Q_{M}}=(\zeta - 1)(-\hat{A}_M+a_{UM}  \hat{p}_U),\\
    \hat{A}_{U,t,\tau}= (\zeta - 1) \frac{\Delta C_{U,t,\tau}}{p_{U,t} Q_{U}}=(\zeta - 1)\hat{p}_U
\end{cases}
\end{equation}
\FL{Sorry I can't follow why we have these RHS in the Eq above.}

\ER{From the dual definition of TFP we get $\hat{A}_M=-\hat{p}_M+\sum_ia_{iM}\hat{p}_i=-\hat{p}_M+a_{UM}\hat{p}_U$, hence, $\hat{p}_M=-\hat{A}_M+a_{UM}\hat{p}_U$. Since $\hat{p}_M=\frac{\Delta C_{M,t,\tau}}{p_{M,t} Q_{M}}$, we have $\frac{\Delta C_{M,t,\tau}}{p_{M,t} Q_{M}}=(-\hat{A}_M+a_{UM}  \hat{p}_U)$.}


\ER{Likewise, $\frac{\Delta C_{U,t,\tau}}{p_{U,t} Q_{U}}=\hat{p}_U$}

As show above, considering the two sectors, we overestimate the aggregate productivity growth
\begin{equation*}
\hat{A}_{t,\tau} =(\zeta - 1)  \frac{\Delta C_{M,t,\tau}+\Delta C_{U,t,\tau}}{\text{GDP}}=(\zeta - 1)  \frac{\Delta \tilde{C}_{M,t,\tau}+2\Delta C_{U,t,\tau}}{\text{GDP}},
\end{equation*}
which is equivalent to $\hat{\tilde{A}}_{t,\tau}= \lambda_M  \hat{A}_M +  2\lambda_U  \hat{A}_U\neq\hat{A}.$

\paragraph{Two solutions}

There are solutions to deal with he double counting, ii) attributing the transition cost to the upstream sectors, which is equivalent to what the classic national accounting does, by removing price change due to the suppliers in the dual TFP formula,
\begin{equation}
\begin{cases}
    \hat{\tilde{A}}_{M,t,\tau} = (\zeta - 1) \frac{\Delta \tilde{C}_{M,t,\tau}}{p_{M,t} Q_{M}}=(\zeta - 1)(\hat{p}_M),\\
    \hat{\tilde{A}}_{U,t,\tau}= (\zeta - 1) \frac{\Delta C_{U,t,\tau}}{p_{U,t} Q_{U}}=(\zeta - 1)\hat{p}_U
\end{cases}
\end{equation}

In this particular case, $\hat{p}_M=-\hat{A}_M$ and $\hat{p}_U=-\hat{A}_U$

\begin{equation*}
 \hat{\tilde{A}}=    (1-\zeta)\left(\lambda_M  \hat{A}_M+  \lambda_U  \hat{A}_U\right),
\end{equation*}
ii) or attributing the cost to the downstream sector, which is equivalent to only considering the downstream sector in the above section,
\begin{equation}
\begin{cases}
    \hat{\tilde{A}}_{M,t,\tau} = (\zeta - 1) \frac{\Delta \tilde{C}_{M,t,\tau}+\Delta C_{U,t,\tau}}{p_{M,t} Q_{M}}=(\zeta - 1)(-\hat{A}_M+a_{UM}  \hat{p}_U)\\
     \hat{\tilde{A}}_{U,t,\tau}=0.
\end{cases}
\end{equation}
Hence, 
\begin{equation*}
 \hat{\tilde{A}}=    \lambda_M  (\zeta - 1)(-\hat{A}_M+a_{UM}  \hat{p}_U)=\lambda_M  (1-\zeta )\hat{A}_M +  \lambda_U  (\zeta - 1)\hat{p}_U=\hat{A}= (1-\zeta)\left(\lambda_M  \hat{A}_M+  \lambda_U  \hat{A}_U\right).
\end{equation*}

% Bibliography ---------------------------

\tiny
\renewcommand\bibname{References}
\bibliographystyle{apalike}
\bibliography{z-bib}
\normalsize
\end{document}
